% archon:covers AlgebraicJacobian/RiemannRoch/SmoothRegular.lean

\chapter{Smooth local algebras and the DVR-stalk substrate}
\label{chap:RiemannRoch_SmoothRegular}

% SOURCE: [Hartshorne], Algebraic Geometry, II.8, Proposition 8.7, p.~173,
% read from references/hartshorne-algebraic-geometry.pdf (PDF page 190 =
% doc p.~173): the conormal isomorphism
% $\mathfrak m/\mathfrak m^2 \cong \Omega_{B/k}\otimes_B k$ at a residue
% field $B/\mathfrak m \cong k$.  The contingency regularity route below
% additionally cites [Stacks Project], Tag 00TT (Algebra,
% `lemma-characterize-smooth-over-field') and Tag 056S (Varieties,
% `lemma-smooth-regular'), read from
% references/stacks-regular-dvr-algebra.tex (lines 38612--38684) and
% references/stacks-smooth-regular-varieties.tex (lines 4596--4614).

\section{Purpose and standing hypotheses}

This chapter isolates the commutative-algebra substrate behind the
codimension-one regularity input of \cref{chap:RiemannRoch_OcOfD}: from the
stalk \(\mathcal O_{C, x}\) of a smooth curve at a codimension-one point it
manufactures a discrete valuation ring. The material is \emph{pure
commutative algebra}: it speaks only of fields, \(k\)-algebras, their modules
of K\"ahler differentials, and the residue dimension of the cotangent space.
No scheme, sheaf, or curve enters here; the stalk-to-substrate bridge is
performed in \cref{chap:RiemannRoch_OcOfD}, not here.

\paragraph{The live route: the conormal isomorphism.} The discrete-valuation
conclusion is reached through the cotangent-space comparison
\[
  \mathfrak m / \mathfrak m^2 \;\cong\; \Omega_{R/k} \otimes_R \kappa
  \qquad (\text{Hartshorne~II.8.7}),
\]
valid at a residue field \(\kappa\) isomorphic to the ground field \(k\) ---
the situation of a \(k\)-rational point over an algebraically closed \(k\).
For \(R\) formally smooth over \(k\) this map is an isomorphism, so the
residue dimension of the cotangent space equals the rank of the locally free
module of differentials:
\[
  \operatorname{finrank}_\kappa
    \bigl(\texttt{IsLocalRing.CotangentSpace}\ R\bigr)
  \;=\;
  \operatorname{finrank}_R \Omega_{R/k}.
\]
When the differentials have rank one --- the relative-dimension-one input
threaded in from the smooth curve --- the cotangent space is one-dimensional,
and the standard cotangent-space characterisation of discrete valuation rings
finishes the argument. This route assembles results already present in
Mathlib's cotangent-complex API; it never forms the Krull dimension and needs
no regularity theory.

\paragraph{A contingency route via regularity.} The closing section records an
independent passage that first proves \emph{regularity} of the local ring by
the Jacobian-criterion argument of Stacks~00TT and only then reads the
cotangent dimension off the identity \(\dim_\kappa \mathfrak m/\mathfrak m^2 =
\operatorname{spanFinrank} \mathfrak m = \operatorname{ringKrullDim} R\). That
route is \emph{not} on the critical path of the headline cone below: it is
superseded by the conormal isomorphism and is retained only as a fallback
(admitting a short proof once the cotangent dimension is known to be one,
since then \(1 = \dim_\kappa \mathfrak m/\mathfrak m^2 =
\operatorname{ringKrullDim} R\) certifies regularity).

\paragraph{Standing hypotheses.} Throughout, \(k\) is a field and \(R\) is a
Noetherian local \(k\)-algebra. The phrase ``\(R\) is smooth over \(k\)''
means that \(R\) is essentially of finite type over \(k\) and formally smooth
over \(k\) --- equivalently, \(R\) is the localisation at a prime of a
finite-type \(k\)-algebra \(S\) that is standard smooth over \(k\) on a basic
open neighbourhood of that prime. We write \(\kappa = R / \mathfrak m\) for
the residue field, \(\Omega_{R/k}\) for the module of K\"ahler differentials,
and \(\operatorname{ringKrullDim} R\) for the Krull dimension. In the
geometric application \(\kappa = k\) is algebraically closed.

\section{The conormal isomorphism over a field-rational residue field}

This section assembles the conormal isomorphism from Mathlib's cotangent
complex. We first anchor the Mathlib ingredients, then state the
isomorphism.

\begin{lemma}[Smooth \(\Rightarrow\) formally smooth]
  \label{lem:smooth_formallySmooth}
  \lean{Algebra.Smooth.formallySmooth}
  \mathlibok
  \textit{Provided by Mathlib
  (\texttt{Mathlib.RingTheory.Smooth.Basic}).}
  If a \(k\)-algebra \(A\) is smooth over \(k\) (\(\texttt{Algebra.Smooth}\
  k\ A\)) then it is formally smooth over \(k\)
  (\(\texttt{Algebra.FormallySmooth}\ k\ A\)). This is the structural
  projection used to feed the formal-smoothness hypothesis of the conormal
  isomorphism.
\end{lemma}

\begin{lemma}[Kernel of the residue map is the maximal ideal]
  \label{lem:ker_residue_mathlib}
  \lean{IsLocalRing.ker_residue}
  \mathlibok
  \textit{Provided by Mathlib
  (\texttt{Mathlib.RingTheory.LocalRing.ResidueField.Basic}).}
  For a local ring \(R\) with maximal ideal \(\mathfrak m\) and residue field
  \(\kappa = R/\mathfrak m\), the kernel of the residue homomorphism
  \(\texttt{residue}\ R : R \to \kappa\) is the maximal ideal:
  \(\operatorname{ker}(\texttt{residue}\ R) = \mathfrak m\). This identifies
  the source of the conormal map (the cotangent module
  \(I/I^2\) of \(I = \operatorname{ker}(R \to \kappa)\)) with
  \(\mathfrak m/\mathfrak m^2\).
\end{lemma}

\begin{lemma}[Formal smoothness of a perfect-field extension]
  \label{lem:formallySmooth_of_perfectField_mathlib}
  \lean{Algebra.FormallySmooth.of_perfectField}
  \mathlibok
  \textit{Provided by Mathlib
  (\texttt{Mathlib.RingTheory.Smooth.Field}).}
  If \(k\) is a perfect field and \(L\) is essentially of finite type over
  \(k\) (\(\texttt{Algebra.EssFiniteType}\ k\ L\)), then \(L\) is formally
  smooth over \(k\) (\(\texttt{Algebra.FormallySmooth}\ k\ L\)). Applied to
  \(L = \kappa\): over an algebraically closed --- hence perfect --- field
  \(k\), the residue field \(\kappa\) (in the application \(\kappa = k\)) is
  formally smooth over \(k\).
\end{lemma}

\begin{lemma}[Formal smoothness kills first cotangent homology]
  \label{lem:subsingleton_h1Cotangent_mathlib}
  \lean{Algebra.FormallySmooth.subsingleton_h1Cotangent}
  \mathlibok
  \textit{Provided by Mathlib
  (\texttt{Mathlib.RingTheory.Smooth.Basic}).}
  If \(A\) is formally smooth over \(R\)
  (\(\texttt{Algebra.FormallySmooth}\ R\ A\)) then the first Andr\'e--Quillen
  (cotangent-complex) homology vanishes:
  \(\texttt{Subsingleton}\ (\texttt{Algebra.H1Cotangent}\ R\ A)\). Applied to
  \(\kappa\) formally smooth over \(k\), this gives
  \(\texttt{Subsingleton}\ (\texttt{H1Cotangent}\ k\ \kappa)\).
\end{lemma}

\begin{lemma}[Conormal map injective iff first cotangent homology vanishes]
  \label{lem:kerCotangentToTensor_injective_iff_mathlib}
  \lean{Algebra.FormallySmooth.kerCotangentToTensor_injective_iff}
  \mathlibok
  \textit{Provided by Mathlib
  (\texttt{Mathlib.RingTheory.Smooth.Basic}).}
  Let \(R \to P \to A\) be \(k\)-algebra maps with \(P\) formally smooth over
  \(R\) and \(\texttt{algebraMap}\ P\ A\) surjective, with kernel \(I\). Then
  the conormal map
  \(\texttt{kerCotangentToTensor}\ R\ P\ A : I/I^2 \to A \otimes_P
  \Omega_{P/R}\) is injective if and only if
  \(\texttt{Subsingleton}\ (\texttt{Algebra.H1Cotangent}\ R\ A)\). Applied
  with \(P = R\), \(A = \kappa\), base \(R = k\), the residue surjection
  \(R \to \kappa\) makes this the injectivity of the conormal map
  \(\mathfrak m/\mathfrak m^2 \to \kappa \otimes_R \Omega_{R/k}\).
\end{lemma}

\begin{lemma}[Right-exact conormal sequence]
  \label{lem:exact_kerCotangentToTensor_mapBaseChange_mathlib}
  \lean{KaehlerDifferential.exact_kerCotangentToTensor_mapBaseChange}
  \mathlibok
  \textit{Provided by Mathlib
  (\texttt{Mathlib.RingTheory.Kaehler.Basic}).}
  For a surjective \(R\)-algebra homomorphism \(A \to B\) with kernel \(I\),
  the conormal sequence
  \[
    I/I^2
      \xrightarrow{\ \texttt{kerCotangentToTensor}\ R\ A\ B\ }
    B \otimes_A \Omega_{A/R}
      \xrightarrow{\ \texttt{mapBaseChange}\ R\ A\ B\ }
    \Omega_{B/R}
  \]
  is exact (the second exact sequence of differentials, Hartshorne~II.8.4A).
  Applied with \(A = R\), \(B = \kappa\), this is the right-exact sequence
  \(\mathfrak m/\mathfrak m^2 \to \kappa \otimes_R \Omega_{R/k} \to
  \Omega_{\kappa/k} \to 0\).
\end{lemma}

\begin{lemma}[Conormal isomorphism
\(\mathfrak m/\mathfrak m^2 \cong \Omega_{R/k}\otimes_R\kappa\)]
  \label{lem:conormal_cotangentSpace_baseChange_kaehler}
  \lean{IsLocalRing.cotangentSpace_linearEquiv_baseChange_kaehler}
  \leanok
  \uses{lem:smooth_formallySmooth, lem:ker_residue_mathlib,
        lem:formallySmooth_of_perfectField_mathlib,
        lem:subsingleton_h1Cotangent_mathlib,
        lem:kerCotangentToTensor_injective_iff_mathlib,
        lem:exact_kerCotangentToTensor_mapBaseChange_mathlib}
  % SOURCE: [Hartshorne], Algebraic Geometry, II.8, Proposition 8.7,
  % p.~173 (read from references/hartshorne-algebraic-geometry.pdf,
  % PDF page 190 = doc p.~173).
  % SOURCE QUOTE: "Let $B$ be a local ring which contains a field $k$
  % isomorphic to its residue field $B/\mathfrak{m}$. Then the map
  % $\delta : \mathfrak{m}/\mathfrak{m}^2 \to \Omega_{B/k} \otimes_B k$ of
  % (8.4A) is an isomorphism."
  \textit{Source: Hartshorne, Algebraic Geometry, II.8.7.}
  Let \(k\) be a field and \(R\) a Noetherian local \(k\)-algebra,
  essentially of finite type and formally smooth over \(k\), with maximal
  ideal \(\mathfrak m\) and residue field \(\kappa\), and suppose \(\kappa\)
  is formally smooth over \(k\) (automatic when \(\kappa = k\), e.g.\ over an
  algebraically closed \(k\) at a \(k\)-rational point). Then the conormal
  map
  \[
    \delta : \mathfrak m / \mathfrak m^2 \longrightarrow
      \Omega_{R/k} \otimes_R \kappa
  \]
  is a \(\kappa\)-linear isomorphism. Equivalently, the cotangent space
  \(\texttt{IsLocalRing.CotangentSpace}\ R = \kappa \otimes_R
  (\mathfrak m/\mathfrak m^2)\) is \(\kappa\)-linearly isomorphic to
  \(\kappa \otimes_R \Omega_{R/k}\).
\end{lemma}

% SOURCE QUOTE PROOF: "According to (8.4A), the cokernel of $\delta$ is
% $\Omega_{k/k} = 0$, so $\delta$ is surjective. To show that $\delta$ is
% injective, it will be sufficient to show that the map
% $\delta' : \operatorname{Hom}_k(\Omega_{B/k} \otimes k, k) \to
% \operatorname{Hom}_k(\mathfrak{m}/\mathfrak{m}^2, k)$ of dual vector spaces
% is surjective. The term on the left is isomorphic to
% $\operatorname{Hom}_B(\Omega_{B/k}, k)$, which by definition of the
% differentials, can be identified with the set $\operatorname{Der}_k(B, k)$
% of $k$-derivations of $B$ to $k$. If $d : B \to k$ is a derivation, then
% $\delta'(d)$ is obtained by restricting to $\mathfrak{m}$, and noting that
% $d(\mathfrak{m}^2) = 0$. Now, to show that $\delta'$ is surjective, let
% $h \in \operatorname{Hom}(\mathfrak{m}/\mathfrak{m}^2, k)$. For any
% $b \in B$, we can write $b = \lambda + c$, $\lambda \in k$,
% $c \in \mathfrak{m}$, in a unique way. Define $db = h(\bar{c})$, where
% $\bar{c} \in \mathfrak{m}/\mathfrak{m}^2$ is the image of $c$. Then one
% verifies immediately that $d$ is a $k$-derivation of $B$ to $k$, and that
% $\delta'(d) = h$. Thus $\delta'$ is surjective, as required."
\begin{proof}
  \leanok
  \uses{lem:smooth_formallySmooth, lem:ker_residue_mathlib,
        lem:formallySmooth_of_perfectField_mathlib,
        lem:subsingleton_h1Cotangent_mathlib,
        lem:kerCotangentToTensor_injective_iff_mathlib,
        lem:exact_kerCotangentToTensor_mapBaseChange_mathlib}
  Realise \(\delta\) as Mathlib's
  \(\texttt{kerCotangentToTensor}\ k\ R\ \kappa\). Its source is
  \(I/I^2\) for \(I = \operatorname{ker}(R \to \kappa)\), and
  \cref{lem:ker_residue_mathlib} identifies \(I = \mathfrak m\), so the source
  is \(\mathfrak m / \mathfrak m^2\) and the target is
  \(\kappa \otimes_R \Omega_{R/k}\), the base change of the differentials.

  \emph{Injectivity.} Smoothness of \(R\) over \(k\) gives formal smoothness
  (\cref{lem:smooth_formallySmooth}). The residue field \(\kappa\), being
  essentially of finite type over the perfect field \(k\), is formally smooth
  over \(k\) (\cref{lem:formallySmooth_of_perfectField_mathlib}; trivially so
  when \(\kappa = k\)), whence the first cotangent homology vanishes,
  \(\texttt{Subsingleton}\ (\texttt{H1Cotangent}\ k\ \kappa)\)
  (\cref{lem:subsingleton_h1Cotangent_mathlib}). By
  \cref{lem:kerCotangentToTensor_injective_iff_mathlib} (applied to the
  surjection \(R \to \kappa\)) the map \(\delta\) is therefore injective. This
  is Hartshorne's dual-derivation argument repackaged through the
  cotangent-complex vanishing.

  \emph{Surjectivity.} The conormal sequence
  \(\mathfrak m/\mathfrak m^2 \to \kappa \otimes_R \Omega_{R/k} \to
  \Omega_{\kappa/k} \to 0\) is exact
  (\cref{lem:exact_kerCotangentToTensor_mapBaseChange_mathlib}). Since
  \(\kappa = k\), the relative differentials \(\Omega_{\kappa/k} =
  \Omega_{k/k}\) vanish, so the base-change map kills the middle term's
  quotient and the range of \(\delta\) is everything; \(\delta\) is
  surjective. (This is the cokernel computation
  \(\operatorname{coker}\delta = \Omega_{k/k} = 0\) of the source proof.)

  Being injective and surjective, \(\delta\) is a \(\kappa\)-linear
  isomorphism. Applying \(\kappa \otimes_R (-)\) transports it to the stated
  isomorphism \(\texttt{CotangentSpace}\ R = \kappa \otimes_R
  (\mathfrak m/\mathfrak m^2) \cong \kappa \otimes_R \Omega_{R/k}\).
\end{proof}

\section{Cotangent dimension equals the K\"ahler rank}

\begin{lemma}[Finite rank under base change]
  \label{lem:finrank_baseChange_mathlib}
  \lean{Module.finrank_baseChange}
  \mathlibok
  \textit{Provided by Mathlib
  (\texttt{Mathlib.LinearAlgebra.Dimension.Constructions}).}
  For a ring homomorphism \(R \to \kappa\) and an \(R\)-module \(M\), the
  \(\kappa\)-rank of the base change equals the \(R\)-rank of \(M\):
  \(\operatorname{finrank}_\kappa (\kappa \otimes_R M) =
  \operatorname{finrank}_R M\).
\end{lemma}

\begin{lemma}[Cotangent dimension equals the differential rank]
  \label{lem:finrank_cotangentSpace_eq_rank_kaehler}
  \lean{IsLocalRing.finrank_cotangentSpace_eq_finrank_kaehler}
  \leanok
  \uses{lem:conormal_cotangentSpace_baseChange_kaehler,
        lem:finrank_baseChange_mathlib}
  \textit{Source: Hartshorne, Algebraic Geometry, II.8.7 (numerical
  corollary).}
  Under the hypotheses of
  \cref{lem:conormal_cotangentSpace_baseChange_kaehler}, the residue-field
  dimension of the cotangent space equals the \(R\)-rank of the
  differentials:
  \[
    \operatorname{finrank}_\kappa
      \bigl(\texttt{IsLocalRing.CotangentSpace}\ R\bigr)
    \;=\;
    \operatorname{finrank}_R \Omega_{R/k}.
  \]
\end{lemma}

\begin{proof}
  \leanok
  \uses{lem:conormal_cotangentSpace_baseChange_kaehler,
        lem:finrank_baseChange_mathlib}
  By \cref{lem:conormal_cotangentSpace_baseChange_kaehler} the cotangent
  space \(\texttt{CotangentSpace}\ R\) is \(\kappa\)-linearly isomorphic to
  \(\kappa \otimes_R \Omega_{R/k}\), so the two have equal
  \(\kappa\)-dimension. The base-change rank identity
  \cref{lem:finrank_baseChange_mathlib} computes
  \(\operatorname{finrank}_\kappa (\kappa \otimes_R \Omega_{R/k}) =
  \operatorname{finrank}_R \Omega_{R/k}\). Composing the two equalities
  yields the claim.
\end{proof}

\section{From smooth and relative dimension one to a discrete valuation ring}

The headline output of the chapter feeds the relative-dimension-one
differential input through the cotangent comparison and finishes with the
standard discrete-valuation characterisation. The rank-one input itself is
supplied by the standard-smooth chart, whose differentials are free of rank
equal to the relative dimension and localise to the differentials of the
stalk.

\begin{lemma}[Standard smooth \(\Rightarrow\) free differentials]
  \label{lem:standardSmooth_kaehler_free}
  \lean{Algebra.IsStandardSmooth.free_kaehlerDifferential}
  \mathlibok
  \textit{Provided by Mathlib
  (\texttt{Mathlib.RingTheory.Smooth.StandardSmooth}).}
  If \(S\) is a standard smooth \(R\)-algebra
  (\(\texttt{Algebra.IsStandardSmooth}\ R\ S\)) then its module of K\"ahler
  differentials \(\Omega_{S/R}\) is a free \(S\)-module
  (\(\texttt{Module.Free}\ S\ \Omega_{S/R}\)).
\end{lemma}

\begin{lemma}[Rank of the differentials of a standard smooth algebra]
  \label{lem:standardSmooth_kaehler_rank}
  \lean{Algebra.IsStandardSmoothOfRelativeDimension.rank_kaehlerDifferential}
  \mathlibok
  \textit{Provided by Mathlib
  (\texttt{Mathlib.RingTheory.Smooth.StandardSmooth}).}
  If \(S\) is a nontrivial standard smooth \(R\)-algebra of relative
  dimension \(n\)
  (\(\texttt{Algebra.IsStandardSmoothOfRelativeDimension}\ n\ R\ S\)) then the
  \(S\)-module rank of its differentials equals \(n\):
  \(\operatorname{rank}_S \Omega_{S/R} = n\). Combined with
  \cref{lem:standardSmooth_kaehler_free}, the differentials are free of rank
  \(n\). For \(n = 1\) this is the relative-dimension-one input feeding the
  cotangent comparison.
\end{lemma}

\begin{lemma}[Differentials localise]
  \label{lem:kaehler_isLocalizedModule_mathlib}
  \lean{KaehlerDifferential.isLocalizedModule}
  \mathlibok
  \textit{Provided by Mathlib
  (\texttt{Mathlib.RingTheory.Kaehler.TensorProduct}).}
  If \(S\) is a \(k\)-algebra and \(R\) is a localisation of \(S\) at a
  submonoid \(p\) (\(\texttt{IsLocalization}\ p\ R\)), then \(\Omega_{R/k}\)
  is the localisation of \(\Omega_{S/k}\) at \(p\) (it is an
  \(\texttt{IsLocalizedModule}\) of the differentials of \(S\)). Hence
  freeness and rank of the K\"ahler differentials transport from a
  finite-type chart \(S\) to its localisation \(R = S_{\mathfrak q}\): if
  \(\Omega_{S/k}\) is free of rank \(n\) then so is \(\Omega_{R/k}\).
\end{lemma}

\begin{lemma}[Smooth local \(k\)-algebra of relative dimension one
\(\Rightarrow\) DVR]
  \label{lem:smooth_stalk_isDVR}
  \lean{Algebra.isDiscreteValuationRing_of_smooth_dim_one}
  \leanok
  \uses{lem:smooth_formallySmooth,
        lem:finrank_cotangentSpace_eq_rank_kaehler,
        lem:finrank_cotangentSpace_eq_one_iff_isDVR,
        lem:standardSmooth_kaehler_free,
        lem:standardSmooth_kaehler_rank,
        lem:kaehler_isLocalizedModule_mathlib}
  Let \(k\) be a field and \(R\) a Noetherian local \(k\)-algebra that is an
  integral domain, smooth over \(k\) (essentially of finite type and formally
  smooth), with residue field \(\kappa\) formally smooth over \(k\)
  (automatic when \(\kappa = k\)), and suppose the differentials have rank
  one: \(\operatorname{finrank}_R \Omega_{R/k} = 1\). Then \(R\) is a discrete
  valuation ring (\(\texttt{IsDiscreteValuationRing}\ R\)).
\end{lemma}

\begin{proof}
  \leanok
  \uses{lem:smooth_formallySmooth,
        lem:finrank_cotangentSpace_eq_rank_kaehler,
        lem:finrank_cotangentSpace_eq_one_iff_isDVR,
        lem:standardSmooth_kaehler_free,
        lem:standardSmooth_kaehler_rank,
        lem:kaehler_isLocalizedModule_mathlib}
  Smoothness of \(R\) over \(k\) gives formal smoothness
  (\cref{lem:smooth_formallySmooth}), so the conormal comparison applies. By
  \cref{lem:finrank_cotangentSpace_eq_rank_kaehler},
  \[
    \operatorname{finrank}_\kappa
      \bigl(\texttt{CotangentSpace}\ R\bigr)
    = \operatorname{finrank}_R \Omega_{R/k} = 1 .
  \]
  The ring \(R\) is a Noetherian local integral domain, so the
  cotangent-space characterisation
  \cref{lem:finrank_cotangentSpace_eq_one_iff_isDVR} (Mathlib's
  \(\texttt{IsLocalRing.finrank\_CotangentSpace\_eq\_one\_iff}\)) converts the
  value \(1\) into \(\texttt{IsDiscreteValuationRing}\ R\).

  \emph{Provenance of the rank-one hypothesis.} In the geometric application
  the input \(\operatorname{finrank}_R \Omega_{R/k} = 1\) is not assumed
  abstractly but produced from the relative-dimension-one standard-smooth
  chart \(S\) with \(R = S_{\mathfrak q}\): the chart differentials
  \(\Omega_{S/k}\) are free of rank one
  (\cref{lem:standardSmooth_kaehler_free},
  \cref{lem:standardSmooth_kaehler_rank}), and they localise to
  \(\Omega_{R/k}\) (\cref{lem:kaehler_isLocalizedModule_mathlib}), so
  \(\Omega_{R/k}\) is free of rank one and
  \(\operatorname{finrank}_R \Omega_{R/k} = 1\). This is the same geometric
  datum the contingency regularity route also consumes, but here it enters as
  a present module fact rather than feeding a regularity machine.
\end{proof}

\section{Contingency: smooth implies regular (off the critical path)}

The remainder of the chapter records the regularity passage. It is
\emph{superseded} by the conormal isomorphism above and is \emph{not} required
by \cref{lem:smooth_stalk_isDVR} or by the headline cone of
\cref{chap:RiemannRoch_OcOfD}; it is retained as a contingency, since
once the cotangent dimension is known to be one the identity
\(\dim_\kappa \mathfrak m/\mathfrak m^2 = \operatorname{ringKrullDim} R\)
certifies regularity in reverse.

\begin{theorem}

\leanok
[Smooth local \(k\)-algebra \(\Rightarrow\) regular local ring]
  \label{thm:smooth_local_isRegularLocalRing}
  \lean{Algebra.IsRegularLocalRing_of_smooth}
  \uses{lem:smooth_formallySmooth, lem:standardSmooth_kaehler_free,
        lem:standardSmooth_kaehler_rank,
        lem:isRegularLocalRing_iff_finrank_cotangentSpace}
  % SOURCE: [Stacks Project], Tag 00TT = Lemma 10.140.3
  % (`algebra-lemma-characterize-smooth-over-field', Section 10.140
  % ``Smooth algebras over fields''). Read from
  % references/stacks-regular-dvr-algebra.tex, lines 38612--38630.
  % SOURCE QUOTE: "Let $k$ be any field.
  % Let $S$ be a finite type $k$-algebra.
  % Let $X = \Spec(S)$.
  % Let $\mathfrak q \subset S$ be a prime
  % corresponding to $x \in X$.
  % The following are equivalent:
  % \begin{enumerate}
  % \item The $k$-algebra $S$ is smooth at $\mathfrak q$ over $k$.
  % \item We have
  % $\dim_{\kappa(\mathfrak q)} \Omega_{S/k} \otimes_S \kappa(\mathfrak q)
  % \leq \dim_x X$.
  % \item We have
  % $\dim_{\kappa(\mathfrak q)} \Omega_{S/k} \otimes_S \kappa(\mathfrak q)
  % = \dim_x X$.
  % \end{enumerate}
  % Moreover, in this case the local ring $S_{\mathfrak q}$ is regular."
  % SOURCE: [Stacks Project], Tag 056S = Lemma 33.25.3
  % (`varieties-lemma-smooth-regular', Section 33.25 ``Schemes smooth over
  % fields''). Read from
  % references/stacks-smooth-regular-varieties.tex, lines 4596--4603.
  % SOURCE QUOTE: "Let $X \to \Spec(k)$ be a smooth morphism where $k$ is a
  % field. Then $X$ is a regular scheme."
  \textit{Source: Stacks Project, Tag 00TT (Lemma 10.140.3, ``Moreover, in
  this case the local ring \(S_{\mathfrak q}\) is regular'') and Tag 056S
  (Lemma 33.25.3, smooth over a field implies regular). Contingency route,
  not on the critical path.}
  Let \(k\) be a field and let \(R\) be a Noetherian local \(k\)-algebra that
  is smooth over \(k\): \(R\) is essentially of finite type over \(k\) and
  formally smooth over \(k\), realised as the localisation \(R = S_{\mathfrak
  q}\) of a finite-type \(k\)-algebra \(S\) at a prime \(\mathfrak q\) on a
  basic open where \(S\) is standard smooth over \(k\). Then \(R\) is a
  regular local ring (\(\texttt{IsRegularLocalRing}\ R\)).
\end{theorem}

% SOURCE QUOTE PROOF: "If $S$ is smooth at $\mathfrak q$ over $k$, then there
% exists a $g \in S$, $g \not \in \mathfrak q$ such that $S_g$ is standard
% smooth over $k$, see Lemma \ref{lemma-smooth-syntomic}. A standard smooth
% algebra over $k$ has a module of differentials which is free of rank equal
% to the dimension, see Lemma \ref{lemma-standard-smooth} (use that a
% relative global complete intersection over a field has dimension equal to
% the number of variables minus the number of equations). Thus we see that
% (1) implies (3). To finish the proof of the lemma it suffices to show that
% (2) implies (1) and that it implies that $S_{\mathfrak q}$ is regular.
% \medskip\noindent
% Assume (2). By Nakayama's Lemma \ref{lemma-NAK} we see that
% $\Omega_{S/k, \mathfrak q}$ can be generated by $\leq \dim_x X$ elements.
% We may replace $S$ by $S_g$ for some $g \in S$, $g \not \in \mathfrak q$
% such that $\Omega_{S/k}$ is generated by at most $\dim_x X$ elements.
% Let $K/k$ be an algebraically closed field extension such that there exists
% a $k$-algebra map $\psi : \kappa(\mathfrak q) \to K$. Consider
% $S_K = K \otimes_k S$. Let $\mathfrak m \subset S_K$ be the maximal ideal
% corresponding to the surjection
% $S_K = K \otimes_k S \to K \otimes_k \kappa(\mathfrak q)
% \xrightarrow{\text{id}_K \otimes \psi} K$.
% Note that $\mathfrak m \cap S = \mathfrak q$, in other words $\mathfrak m$
% lies over $\mathfrak q$. By Lemma
% \ref{lemma-dimension-at-a-point-preserved-field-extension} the dimension of
% $X_K = \Spec(S_K)$ at the point corresponding to $\mathfrak m$ is
% $\dim_x X$. By Lemma \ref{lemma-dimension-closed-point-finite-type-field}
% this is equal to $\dim((S_K)_{\mathfrak m})$. By Lemma
% \ref{lemma-differentials-base-change} the module of differentials of $S_K$
% over $K$ is the base change of $\Omega_{S/k}$, hence also generated by at
% most $\dim_x X = \dim((S_K)_{\mathfrak m})$ elements. By Lemma
% \ref{lemma-characterize-smooth-kbar} we see that $S_K$ is smooth at
% $\mathfrak m$ over $K$. By Lemma \ref{lemma-flat-base-change-locus-smooth}
% this implies that $S$ is smooth at $\mathfrak q$ over $k$. This proves (1).
% Moreover, we know by Lemma \ref{lemma-characterize-smooth-kbar} that the
% local ring $(S_K)_{\mathfrak m}$ is regular. Since
% $S_{\mathfrak q} \to (S_K)_{\mathfrak m}$ is flat we conclude from Lemma
% \ref{lemma-flat-under-regular} that $S_{\mathfrak q}$ is regular."
\begin{proof}
  \uses{lem:smooth_formallySmooth, lem:standardSmooth_kaehler_free,
        lem:standardSmooth_kaehler_rank,
        lem:isRegularLocalRing_iff_finrank_cotangentSpace}
  We follow Stacks~00TT. Smoothness of \(R = S_{\mathfrak q}\) over \(k\)
  unfolds (\cref{lem:smooth_formallySmooth}) to a basic open on which the
  finite-type chart \(S_g\) (\(g \notin \mathfrak q\)) is standard smooth over
  \(k\); replacing \(S\) by \(S_g\), we may assume \(S\) itself is standard
  smooth of some relative dimension \(n\).

  \emph{Step 1 (differentials of the chart).} A standard smooth \(k\)-algebra
  has a free module of K\"ahler differentials
  (\cref{lem:standardSmooth_kaehler_free}), and when it is standard smooth of
  relative dimension \(n\) the rank is exactly \(n\)
  (\cref{lem:standardSmooth_kaehler_rank}). Thus \(\Omega_{S/k}\) is free of
  rank \(n\), and \(\dim_{\kappa(\mathfrak q)} \Omega_{S/k} \otimes_S
  \kappa(\mathfrak q) = n = \dim_x X\); this is the equality~(3) of the Stacks
  quotation, and supplies the numerical input of the Jacobian criterion.

  \emph{Step 2 (regularity at the prime, Stacks~00TT).} Stacks~00TT then
  concludes that the local ring \(S_{\mathfrak q}\) is regular. Its argument
  base-changes to an algebraically closed extension \(K / k\) with a
  \(k\)-embedding \(\psi : \kappa(\mathfrak q) \to K\): at the maximal ideal
  \(\mathfrak m \subset S_K = K \otimes_k S\) lying over \(\mathfrak q\), the
  ring \(S_K\) is smooth over the algebraically closed field \(K\), so the
  local ring \((S_K)_{\mathfrak m}\) is regular (smoothness over an
  algebraically closed field at a closed point gives regularity, via the
  Jacobian rank computation). Since \(S_{\mathfrak q} \to (S_K)_{\mathfrak m}\)
  is faithfully flat, regularity descends to \(S_{\mathfrak q} = R\) (flat
  local homomorphism with regular target and regular fibre forces a regular
  source). This is precisely Tag~056S read pointwise: smooth over a field
  forces every local ring to be regular.

  \emph{Sub-gaps to decompose.} Two ingredients of Step~2 are genuine
  commutative-algebra obligations that do not follow immediately from standard
  results: (i) the
  regularity of a local ring of a $K$-algebra smooth over an algebraically
  closed field $K$ at a closed point (the Jacobian-criterion regularity,
  Stacks~00TT, the closed-point smoothness-over-$\bar K$ case); and (ii) the descent of
  regularity along the faithfully flat local map $S_{\mathfrak q} \to
  (S_K)_{\mathfrak m}$ (flat descent of regularity, Stacks~00TT), together with the
  transport of regularity from the finite-type chart to its localisation at
  $\mathfrak q$. Hartshorne records the localisation half of (ii) as
  Theorem~II.8.14A. These are the real sub-obligations of this theorem; the
  remaining numerical content (Steps~1) is supplied by the standard results
  below. Because the headline cone runs through the conormal isomorphism,
  these sub-obligations need not be discharged for the discrete-valuation
  conclusion.
\end{proof}

\subsection*{Cotangent dimension of a regular local ring}

The regularity route reads the cotangent dimension off regularity, with no
differential-theoretic input. It rests on two Mathlib facts: for a regular
local ring the minimal number of generators of the maximal ideal equals the
Krull dimension, and for any Noetherian local ring that minimal number equals
the residue dimension of the cotangent space.

\begin{lemma}[Regular local ring: spanFinrank of \(\mathfrak m\) equals
Krull dimension]
  \label{lem:isRegularLocalRing_spanFinrank_eq_ringKrullDim}
  \lean{IsRegularLocalRing.spanFinrank_maximalIdeal}
  \mathlibok
  \textit{Provided by Mathlib
  (\texttt{Mathlib.RingTheory.RegularLocalRing.Defs}).}
  For a regular local ring \(R\) the minimal number of generators of the
  maximal ideal equals the Krull dimension:
  \(\operatorname{spanFinrank}(\mathfrak m_R) =
  \operatorname{ringKrullDim} R\) (as elements of \(\mathbb N \cup
  \{\infty\}\)). This is, in effect, the defining numerical property of a
  regular local ring.
\end{lemma}

\begin{lemma}[spanFinrank of \(\mathfrak m\) equals cotangent dimension]
  \label{lem:spanFinrank_eq_finrank_cotangentSpace}
  \lean{IsLocalRing.spanFinrank_maximalIdeal_eq_finrank_cotangentSpace}
  \mathlibok
  \textit{Provided by Mathlib
  (\texttt{Mathlib.RingTheory.Nakayama} /
  \texttt{Mathlib.RingTheory.LocalRing.CotangentSpace}).}
  For a Noetherian local ring \(R\) the minimal number of generators of the
  maximal ideal equals the residue-field dimension of the cotangent space:
  \(\operatorname{spanFinrank}(\mathfrak m_R) =
  \operatorname{finrank}_{\kappa}
  \bigl(\texttt{IsLocalRing.CotangentSpace}\ R\bigr)\), the latter being
  \(\dim_\kappa \mathfrak m_R / \mathfrak m_R^2\). This is Nakayama's lemma in
  its minimal-generators form.
\end{lemma}

\begin{lemma}

\leanok
[Cotangent dimension of a regular local ring equals its Krull dimension]
  \label{lem:finrank_cotangent_eq_ringKrullDim_of_regular}
  \lean{IsRegularLocalRing.finrank_cotangentSpace_eq_ringKrullDim}
  \uses{lem:isRegularLocalRing_spanFinrank_eq_ringKrullDim,
        lem:spanFinrank_eq_finrank_cotangentSpace}
  Let \(R\) be a regular local ring. Then the residue-field dimension of the
  cotangent space equals the Krull dimension:
  \[
    \operatorname{finrank}_{\kappa}
      \bigl(\texttt{IsLocalRing.CotangentSpace}\ R\bigr)
    \;=\;
    \operatorname{ringKrullDim} R .
  \]
\end{lemma}

\begin{proof}
  \uses{lem:isRegularLocalRing_spanFinrank_eq_ringKrullDim,
        lem:spanFinrank_eq_finrank_cotangentSpace}
        \leanok
  A regular local ring is Noetherian, so
  \cref{lem:spanFinrank_eq_finrank_cotangentSpace} gives
  \(\operatorname{spanFinrank}(\mathfrak m_R) = \operatorname{finrank}_\kappa
  (\texttt{CotangentSpace}\ R)\), and
  \cref{lem:isRegularLocalRing_spanFinrank_eq_ringKrullDim} gives
  \(\operatorname{spanFinrank}(\mathfrak m_R) = \operatorname{ringKrullDim}
  R\). Transitivity of equality yields the claim. (Mathlib also packages this
  composite directly as the forward direction of
  \cref{lem:isRegularLocalRing_iff_finrank_cotangentSpace}.)
\end{proof}

\begin{lemma}[Regularity criterion via cotangent dimension]
  \label{lem:isRegularLocalRing_iff_finrank_cotangentSpace}
  \lean{IsRegularLocalRing.iff_finrank_cotangentSpace}
  \mathlibok
  \textit{Provided by Mathlib
  (\texttt{Mathlib.RingTheory.RegularLocalRing.Defs}).}
  For a Noetherian local ring \(R\), regularity is equivalent to the
  cotangent dimension matching the Krull dimension:
  \(\texttt{IsRegularLocalRing}\ R \iff \operatorname{finrank}_\kappa
  (\texttt{IsLocalRing.CotangentSpace}\ R) = \operatorname{ringKrullDim} R\).
  The forward direction is exactly
  \cref{lem:finrank_cotangent_eq_ringKrullDim_of_regular}; the reverse
  direction supplies the regularity certificate consumed in Step~2 of
  \cref{thm:smooth_local_isRegularLocalRing}, and is also the cheap closure of
  that theorem once the cotangent dimension is pinned to one by the conormal
  route.
\end{lemma}

\subsection*{A normality fallback}

If the Jacobian-criterion regularity of
\cref{thm:smooth_local_isRegularLocalRing} proves harder to formalise than
the discrete-valuation-ring conclusion of \cref{lem:smooth_stalk_isDVR}
requires, the dimension-one case admits a further alternative that never
mentions regularity. A \(k\)-algebra smooth over a field is geometrically
normal, hence (being a local domain) integrally closed in its fraction
field; a Noetherian local integrally closed domain possessing a unique
nonzero prime --- automatic in Krull dimension one --- is a discrete
valuation ring. In Mathlib this is the clause
``\(\texttt{IsIntegrallyClosed}\ R \wedge \exists!\, P,\ P \neq \bot \wedge
P.\texttt{IsPrime}\)'' of the discrete-valuation-ring \texttt{TFAE}
(\cref{lem:isDiscreteValuationRing_TFAE}). This route is recorded as a further
contingency; the live target is \cref{lem:smooth_stalk_isDVR} via the conormal
isomorphism.

\begin{lemma}[Discrete valuation ring characterisations]
  \label{lem:isDiscreteValuationRing_TFAE}
  \lean{IsDiscreteValuationRing.TFAE}
  \mathlibok
  \textit{Provided by Mathlib
  (\texttt{Mathlib.RingTheory.DiscreteValuationRing.TFAE}).}
  For a Noetherian local integral domain \(R\) that is not a field, the
  following are equivalent: \(R\) is a discrete valuation ring; \(R\) is a
  valuation ring; \(R\) is a Dedekind domain; \(R\) is integrally closed with
  a unique nonzero prime ideal; the maximal ideal of \(R\) is principal; the
  cotangent space \(\mathfrak m / \mathfrak m^2\) is one-dimensional over the
  residue field; every nonzero ideal of \(R\) is a power of the maximal
  ideal. The integrally-closed clause is the basis of the normality fallback
  route above.
\end{lemma}
